\section{Results}

Our benchmark results offer a comparative view of summarization performance across all evaluated models. We first examined the correlations among the chosen evaluation metrics and, based on the findings, identified the best-performing model across lexical, semantic, and factual dimensions. Next, we compared model performance across categories and families to identify significant differences. We additionally performed a data leakage analysis to ensure that inclusion of benchmarking data in training data had no impact on the results. We then analyzed how well automatic metrics aligned with human judgment based on expert level result assessment, and finally assessed whether an \ac{LLM}-as-a-judge panel could reproduce these rankings.

\subsection{Metric Correlations}

Pairwise Spearman correlations were computed to analyze the relationships between the different evaluation metrics (Figure~\ref{fig:metric_correlation}). Strong positive correlations were observed among most lexical-based metrics (\ac{ROUGE}-1/2/L, \ac{METEOR}, and \ac{BLEU}), with the correlation between \ac{METEOR} and \ac{BLEU} marking an exception ($\rho = 0.23$). \ac{ROUGE} variants showed almost identical behavior ($\rho \geq 0.88$), while \ac{BLEU} and \ac{METEOR} demonstrated slightly weaker but still substantial alignment with \ac{ROUGE} measures ($\rho = 0.51 - 0.82$).

Most semantic-based metrics (\ac{RoBERTa}, \ac{DeBERTa}, and all-mpnet-base-v2) showed high internal consistency ($\rho > 0.5$), reflecting their shared focus on semantic similarity. When compared with lexical-based metrics, correlations were moderate to strong in most cases, indicating that both categories capture related but not identical dimensions of summary quality. 

Concerning factual consistency metrics, AlignScore and the two MiniCheck variants exhibited very strong correlations with each other ($\rho = 0.82 - 0.94$), while SummaC showed moderately positive correlations with the other factual metrics ($\rho = 0.41 -0.69$). This likely reflects methodological differences, as SummaC relies on zero-shot natural language inference over sentence pairs, while MiniCheck and AlignScore use models specifically optimized for factual consistency assessment. When compared with lexical- and semantic-based metrics, correlations were generally weaker and, in some cases, even slightly negative. This pattern can largely be attributed to the different point of reference used by factual metrics.

Overall, these relationships demonstrate that the various metrics are broadly consistent while providing complementary perspectives. This supports the use of an aggregated ``Overall Performance Score'' as a balanced indicator of overall summarization performance.

\subsection{Overall Model Performance}

Based on overall performance ranks, referred to as ``Performance Rank" and derived from our multi-metric evaluation framework (\ref{subsec:overall_performance_metric} Overall Performance Metric), models from the Mistral family were top-ranked across the three evaluated metric dimensions, as depicted in Figure~\ref{fig:rank_heatmap}. 

The mistral-small3.2:24b model ranked first, followed by mistral-small-2506 and mistral-medium-2505. The lowest-ranked models included OpenBioLLM-Llama3-8B, bigbird-pegasus-large-pubmed and pegasus-pubmed from the \ac{PEGASUS} family, and mT5\_multilingual\_XLSum from the \ac{T5} family. Overall, domain-specific \acp{SLM} and \ac{EDM} models showed poor performance across all the different dimensions of metrics.

Among the ten top-ranked models, six were general-purpose \acp{LLM} and four were general-purpose \acp{SLM}. A similar trend was observed in the top half of the ranking (positions 1 to 31) with the presence of just a few reasoning-oriented and domain-specific models. In contrast, in the lower half of the ranking, where models performed poorly across most metric dimensions, the majority were reasoning-oriented \acp{SLM}, general-purpose \acp{EDM}, domain-specific \acp{EDM}, and traditional models. 

Considering each metric dimension separately, several Mistral models ranked high in the lexical dimension, with GPT-5-nano, a reasoning-oriented \ac{LLM} from the \ac{GPT} family, ranking third. In the semantic dimension, Mistral models achieved high ranks, but the highest positions were occupied by the Phi4:14b general-purpose \ac{LLM} and several Granite models, including granite4:tiny-h and granite3.3:2b. Finally, the highest ranks based on the factual dimension were covered by traditional models and some \acp{EDM}. Detailed performance indications for each model across all individual ranked metrics are provided in Supplementary Figure~\ref{fig:supplementary_overview}.

Models could also decline to summarize by returning the INSUFFICIENT\_FINDINGS token. This was rare overall, occurring in 736 of 62,000 generated summaries (1.2\%) and exclusively among prompt-based models. Usage was concentrated in the two domain-specific biomedical models OpenBioLLM-Llama3-8B (20.0\%) and medllama2 (9.7\%), while all other models returned it in fewer than 6\% of cases.

\subsection{Category Comparisons}

To compare model performance across categories, we displayed the distribution of overall performance scores within each category, along with the best- and worst-performing individual models (Figure~\ref{fig:category_boxplots_and_gameshowell}a). Results of the Games-Howell post-hoc comparisons between individual groups are given in Figure~\ref{fig:category_boxplots_and_gameshowell}b.

All general-purpose \acp{LLM} achieved high overall scores, thus forming the top category of models, with general-purpose \acp{SLM} following second, however showing a slightly wider spread in performance scores. Domain-specific \acp{LLM} and reasoning-oriented \acp{LLM} followed in third and fourth place based on average overall performance.

Domain-specific \acp{EDM}, domain-specific \acp{SLM}, traditional models, and general-purpose \acp{EDM} performed significantly worse as compared to general-purpose \acp{LM}. Domain-specific \acp{EDM} and domain-specific \acp{SLM} displayed the widest spread of performance (Figure~\ref{fig:category_boxplots_and_gameshowell}a).

\subsection{Family Comparisons}

To complement the category-level findings, we examined performance across model families, since models within the same family often share architectural features or training strategies that may cause consistent performance trends. Figure~\ref{fig:family_boxplots_and_gameshowell}a shows the distribution of overall performance scores across all model families. Similar to the category-level analysis, there were clear performance differences between architectural lineages. The Mistral family achieved the strongest overall results, with several members ranking among the top performers in the benchmark. Families dominated by modern \ac{LLM} or \ac{SLM} architectures, such as Granite, Claude, and Gemma, consistently outperformed more traditional extractive approaches and families built on encoder-decoder architectures such as \ac{LED}, \ac{BART}, \ac{T5} and \ac{PEGASUS}.

In particular, the \ac{GPT} and \ac{Llama} families showed lower performance scores than expected given the competitive performance of their top models. This was primarily due to the inclusion of domain-specific small models, such as BioGPT and OpenBioLLM-Llama3-8B, which performed substantially worse than their general-purpose counterparts and therefore pulled down the family-level averages.

Regarding the Games-Howell post-hoc matrix in the family comparison (Figure~\ref{fig:family_boxplots_and_gameshowell}b), a combination of significant and non-significant differences between families was observed, reflecting the greater heterogeneity within some lineages. While many high-performing families differed significantly from traditional and encoder-decoder dominated families, several comparisons among modern \acp{SLM} and \acp{LLM} dominated families did not reach statistical significance.

\subsection{Data Leakage Assessment}

\subsubsection{Cutoff Overlap Results}

Across the evaluated models, potential overlap ranged from minimal levels of 5\% for earlier models with training cutoffs prior to 2024, up to 40.2\% for more recent \acp{LLM} with reported cutoff dates in early 2025. Most models showed low overlap ($\leq 20\%$), indicating that the majority of benchmark articles were published after their respective training cutoff dates. A subset of eight models, including recent versions of Claude, Deepseek, and one \ac{GPT} model, exhibited moderate levels of potential overlap ($24 - 40.2\%$) and were further analyzed via an abstract completion probe in the next step.

A complete overview of training cutoff dates and corresponding potential overlap values for all models is provided in Supplementary Table~\ref{tab:training_cutoff_overlap}.

\subsubsection{Abstract Completion Probe}

To assess whether high potential overlap translated into memorization effects, we conducted an abstract completion probe for the eight models showing moderate levels of potential overlap. For each model \ac{ROUGE}-L scores between generated completions and the original abstracts were compared for articles published before and after the respective training cutoff dates.

Average \ac{ROUGE}-L scores were nearly identical for pre-cutoff and post-cutoff articles (both at approximately 13.5\%) and there were only negligible differences at the individual model level (Table~\ref{tab:completion_probe}). Statistical testing using a one-sided Mann-Whitney U test revealed no significant increase in similarity for pre-cutoff articles in any of the tested models. Effect sizes were consistently close to zero, which further indicates that there is no systematic difference between the two groups.

If memorization effects had occurred, substantially higher similarity scores would be expected for abstracts from pre-cutoff articles. As this was not the case, we concluded that the models were not reproducing memorized content from their training data.

\subsection{Expert Assessment Results}

The results of the expert assessment are summarized in Figure~\ref{fig:assessment_results} and Table~\ref{tab:human_eval}. Values are reported as mean $\pm$ \ac{SEM}. The two Mistral-based models (mistral-small3.2:24b and mistral-small-2506) consistently outperformed the two \acp{EDM} (bigbird-pegasus-large-pubmed and mT5\_multilingual\_XLSum) across all four evaluation dimensions.

For coherence and fluency, both Mistral models scored substantially higher than the \acp{EDM}. In coherence, they reached scores of 4.83 $\pm$ 0.03 and 4.82 $\pm$ 0.04, compared with 3.20 $\pm$ 0.10 for mT5 and 2.08 $\pm$ 0.08 for bigbird-pegasus. The same pattern was observed for fluency, where the Mistral models achieved scores of 4.86 $\pm$ 0.03 and 4.90 $\pm$ 0.03, clearly exceeding mT5 (3.67 $\pm$ 0.09) and bigbird-pegasus (2.23 $\pm$ 0.07). All pairwise comparisons between Mistral models and \acp{EDM} were statistically significant (all $p < 0.001$ after Bonferroni correction). The two \acp{EDM} also differed significantly from each other in both dimensions, while no significant differences were found between the two Mistral models.

Differences were even more pronounced for relevance and consistency. The Mistral models again achieved high scores in relevance (4.07 $\pm$ 0.07 and 4.09 $\pm$ 0.07), while mT5 (1.54 $\pm$ 0.06) and bigbird-pegasus (1.40 $\pm$ 0.05) scored markedly lower. A similar pattern was seen for consistency, where the Mistral models reached scores of 4.56 $\pm$ 0.06 and 4.59 $\pm$ 0.05, compared to 2.17 $\pm$ 0.09 for mT5 and 1.41 $\pm$ 0.05 for bigbird-pegasus. Again, all Mistral versus \ac{EDM} comparisons were statistically significant ($p < 0.001$). The \acp{EDM} differed significantly in consistency, but not in relevance ($p = 0.7308$), and no significant differences were found between the two Mistral models.

Inter-rater agreement was strong on both measures. Gwet's AC2 indicated substantial absolute agreement across all dimensions, ranging from 0.713 (relevance) to 0.794 (fluency), with an overall value of AC2 = 0.725. Mean pairwise Spearman correlations were likewise strong across all dimensions, ranging from $\rho = 0.707$ (coherence) to $\rho = 0.801$ (consistency), with an overall mean of $\rho = 0.765$ (all $p < 0.001$). These results show that evaluators consistently agreed both on the absolute quality of the summaries and on the relative ranking between models.

\subsection{LLM-as-a-Judge Evaluation}
The LLM-as-a-judge panel agreed closely with expert assessments on every dimension (Table~\ref{tab:judge-expert}). To benchmark against human reliability, we computed a leave-one-out ceiling on the same target the judges were scored against: each expert's ratings were correlated against the mean of the remaining experts. The judge's mean agreement ($\rho=0.87$) exceeded this ceiling ($\rho=0.82$), leading to the assumption that it tracked the expert consensus at least as closely as a typical individual expert.


The judge panel did not separate the four dimensions well. Its agreement with the matching expert dimension (diagonal $\rho=0.87$) was barely higher than its agreement with the other dimensions (off-diagonal $\rho=0.85$), a gap of only $0.02$. In effect the judge measured a single overall-quality signal rather than four independent constructs, so its per-dimension scores should not be over-interpreted.

We questioned whether the LLM-as-a-judge panel was capable of ranking the 62 models in a comparable manner to the evaluation metrics (Table~\ref{tab:concordance}). The judge ranking agreed strongly with the lexical, semantic, and composite performance ranks ($\rho=0.75$-$0.83$). The factual metrics were an exception ($\rho=0.09$, n.s.), as they reward verbatim copying, so extractive systems (TextRank, frequency, T5, PEGASUS) ranked highly on factuality but performed poorly under the judge - a construct difference rather than a judge failure.

To analyze the tendency of judges to favor their own model family, we computed, for each judge, how far its score sat above or below the mean of the other two judges, and contrasted its own-family models with all others (Table~\ref{tab:selfpref}). Anthropic exhibited a clear self-preference effect ($+0.30$ on the 1-5 scale), whereas OpenAI and Mistral showed no detectable differential effect ($\leq 0.08$ in magnitude). The direction is consistent with self-preference reported for rubric-based judging \citep{panicksseryLLMEvaluatorsRecognize2024,pombalSelfPreferenceBiasRubricBased2026}. Taking the median across the three judges damped this bias in the final scores.

